\documentclass{article}

\usepackage{calculator, calculus}
\usepackage{amsmath, ifthen, xcolor, titling}

\title{\bf Arithmetic Operations and Scientific 				Calculations}
\date{}

\begin{filecontents}{other_ops.tex}
\begin{itemize}
	\item \verb|\SUBTRACTfunction{\f1}{\f2}{\newSubFun}|
	$F(t)= f1 - f2$
	\item \verb|\PRODUCTfunction{\f1}{\f2}{\newProdFun}|
	$F(t)= f1 * f2$
	\item \verb|\SCALEfunction{num}{\func}{\ScFunc}|
	$F(t) = num * func$
\end{itemize}
\end{filecontents}

\setlength\droptitle{-3cm}

\begin{document}
	\maketitle
	
	\section{The \texttt{calculator} Package}
	
	\subsection{Basic Arithmetic}
	\begin{description}
		\item[Storing Numbers]
			We copy 0.5 to \verb|\half| 
			\COPY{0.5}{\half} 
			So, \half \\
			\COPY{\half}{\newhalf} New half is \newhalf
		%%
		\item[Addition] 
			Add (8, 7) to get 
			\ADD{8}{7}{\added} \added
		%%
		\item[Subtraction] 
			Sub (8 from 10) to get
			\SUBTRACT{10}{8}{\subt} \subt  
		%%
		\item[Multiplication] 
			Multiply (7.5, 6.3)	
			to get \MULTIPLY{7.5}{6.3}{\mult} \mult
		%%
		\item[Division] 
			Divide (80, by 7) to get
			\DIVIDE{80}{7}{\div} \div
		%%%
		\item[Maximum] The Max of (7, 9) is 
			\MAX{7}{9}{\nine}. So, \nine
		\item[Minimum] The Min of (5, 3) is
			\MIN{5}{3}{\three}. So, \three
		%%%
		\item[Simplifying Fractions]
			\FRACTIONSIMPLIFY{10}{16}{\numer}{\denom}
			The numerator and denominator are
			(\numer, \denom) 
		%%%
		\item[Result as a New Macro]
			\newcommand{\ValueAdded}{}
			\ADD{8}{8}{\ValueAdded} 
			My new Macro now holds \ValueAdded
	\end{description}
	
	\subsection{Using in New Commands}
	\newcommand{\decision}[3]{
		\POWER{#1}{#2}{\res}
		\COPY{#3}{\thresh}
		\ifthenelse{\res < \thresh}
			{You still in bound, $\res < \thresh$}						{You're Free,  $\res > \thresh$}
	} 
	\decision{1}{3}{50} \\[3pt]
	\decision{3}{3}{25}
	
	\subsection{The Scope of Result Variables}
	\COPY{8}{\eight} 
	\verb|\eight| is \eight 
	
	{
		At First \verb|\eight| is \eight 
		
		\color{red!70}
		%% Changing the same \eight by adding one 
		\ADD{1}{\eight}{\eight} 
		
		After adding one: 											\verb|\eight| becomes \eight\,!!	
	} 
	
	\noindent
	\verb|\eight| is \eight
	
%	{
%		\POWER{3}{5}{\powthreefive} \powthreefive
%	}
%	\powthreefive

	\subsection{Working with Lengths}
	\newlength{\height}
	
	\begin{description}
		\item[Length Addition]
			\LENGTHADD{1in}{50pt}{\height}
			The length \texttt{height} is \the\height
		%
		\item[Length Subtraction]
			\LENGTHSUBTRACT{2cm}{20pt}{\height}
			Now, \texttt{height} is \the\height
		%
		\item[Length Division]	
			\LENGTHDIVIDE{1in}{1cm}{\lenres} %1in=2.54cm
			One Inch is equal to \lenres 
	\end{description}
	
	\pagebreak
	
	\subsection{Operations on Vectors}
	
	\begin{description}
		\item[Assigning 2-D Vectors]
			\VECTORCOPY(3,4)(\vecta, \vectb) 
			My new vector is $(\vecta, \vectb)$ 
		%
		\item[Assigning 3-D Vectors]
			\VECTORCOPY(3,4,2)(\vecta, \vectb, \vectc)
			Making it 3-D $(\vecta, \vectb, \vectc)$
		%
		\item[Vector Addition]
			Adding the new vector to itself. 
			\VECTORADD(\vecta, \vectb)(\vecta, \vectb)%
			(\vecta, \vectb) So, $(\vecta, \vectb)$
		%
		\item[Scalar-Vector Product]
			\SCALARVECTORPRODUCT{2}(3,4)(\vectx, \vecty)
			So, it becomes $(\vectx, \vecty)$ 
	\end{description}
	
	\subsubsection{More than Three Elements}
	\VECTORADD(1,2,3,4)(1,2,3,4)(\von, \vtw, \vth, \vf)
	Adding $(1,2,3,4)$ to itself gives 
		$(\von, \vtw, \vth, \vf)$
		
	\newcommand{\AddFourDimVect}[4]{
		\ADD{#1}{#1}{\resa} \ADD{#2}{#2}{\resb}
		\ADD{#3}{#3}{\resc} \ADD{#4}{#4}{\resd}
		$(\resa, \resb, \resc, \resd)$
	}
	\noindent \\
	While the result must be: 
	\AddFourDimVect{1}{2}{3}{4} 
	
	\subsection{Operations on Matrices}
	
	\begin{description}
		\item[Assigning 2\,x\,2 Matrix]
			\MATRIXCOPY(1,2;3,4)(\a,\b;\c,\d)
			$\begin{bmatrix}
				\a & \b \\
				\c & \d
			\end{bmatrix}$
		%
		\item[Assigning 3\,x\,3 Matrix]
			\MATRIXCOPY(2, 	1, 4;
						3, -5, 5;
						2, -7, 4)%
						(\ela ,\elb ,\elc;
						 \elld,\elle,\ellf;
						 \ellg,\ellh,\elli)
			$\begin{bmatrix}
				\ela  & \elb  & \elc  \\
				\elld & \elle & \ellf \\
				\ellg & \ellh & \elli
			\end{bmatrix}$
		%
		\item[Transpose a Matrix]
			\TRANSPOSEMATRIX(\a,\b;\c,\d)%
							(\at,\bt;\ct,\dt)
			$\begin{bmatrix}
				\at & \bt \\
				\ct & \dt
		 	\end{bmatrix}$
		 %
		 \item[Matrix Addition]
		 	\MATRIXADD(\a,\b;\c,\d)(\a,\b;\c,\d)%
		 		(\a,\b;\c,\d)
		 	$\begin{bmatrix}
				\a & \b \\
				\c & \d
			 \end{bmatrix}$
		%
		\item[Multiply by Scalar]
			\SCALARMATRIXPRODUCT{2}(1,2;3,4)%
					(\a,\b;\c,\d)
			$\begin{bmatrix}
				\a & \b \\
				\c & \d
			 \end{bmatrix}$
	\end{description}
	
	\pagebreak
	
	\section{The calculus Package}
	
	\subsection{Built-in Functions}
	\begin{description}
		\item[Square Function]
			If we have an input $t=4$ and $f(t)=t^2$, 
			and it's derivative is $f'(t)=2t$, then:
			\SQUAREfunction{4}{\fsol}{\Dfsol}
			$f(4)=\fsol$, and $f'(4)=\Dfsol$.
		%
		\newcommand{\two}{2}
		\item[Cube Function]	
			If we have an input $t=2$ and $f(t)=t^3$, 
			and it's derivative is $f'(t)=3t^2$, then:
			\CUBEfunction{\two}{\twosol}{\Dtwosol}
			$f(2)=\twosol$, and $f'(2)=\Dtwosol$.
		%
		\MULTIPLY{\twosol}{2}{\sqrtin}
		\item[Square Root]
			If we have an input $t=\sqrtin$ 
			and $f(t)=\sqrt{t}$, and it's derivative is
			$f'(t)=\frac{1}{2*\sqrt{t}}$ , then:
			\SQRTfunction{\sqrtin}{\sqrtsol}{\Dsqrtsol}
			$f(\sqrtin)=\sqrtsol$, 
			and $f'(\sqrtin)=\Dsqrtsol$
	\end{description}
	
	\subsection{Defining New Functions via Operations}
	\begin{description}
		\item[Summing Two Functions]
			We can define $F(t)=\sqrt{t} + \sin t$ using 
			the \verb|\SUMfunction|.
		   \SUMfunction{\SQRTfunction}{\SINfunction}{\F}
		   \F{4}{\Fsol}{\DFsol}	
		   %
		   So that we have: 
		   \begin{itemize}
		   		\item $F(4) = \sqrt{4} + \sin 4$ 
		   			equal to ~\Fsol
		   		\item $F'(4)= \frac{1}{2*\sqrt{4}} + 							\cos 4$ equal to \DFsol
		   \end{itemize}
		%%
		\item[Some Other Operations] that work the same
		 	way:
		 	%% LaTeX2e file `other_ops.tex'
%% generated by the `filecontents' environment
%% from source `calculator' on 2023/12/28.
%%
\begin{itemize}
 \item \verb|\SUBTRACTfunction{\f1}{\f2}{\newSubFun}|
 $F(t)= f1 - f2$
 \item \verb|\PRODUCTfunction{\f1}{\f2}{\newProdFun}|
 $F(t)= f1 * f2$
 \item \verb|\SCALEfunction{num}{\func}{\ScFunc}|
 $F(t) = num * func$
\end{itemize}

	\end{description}
	
	\subsection{Example: Linear Combination}
	The linear combination of two functions $f(x)$ and
	 $g(x)$ is given by: \\
	 \[ h(x) = a \cdot f(x) + b \cdot g(x) \]
	%
	 \noindent\\[5pt]
	 Compose and Evaluate:
	 $F(t) = 2e^{cos t} - e^{-t}sin t$ 
	%
	%%% 1. First Part: e^(cos t)
	\COMPOSITIONfunction{\EXPfunction}{\COSfunction}				{\FunOne}
	
	%%% 2. Secoond Part: e^(-t) * sin(t)
	\SCALEVARIABLEfunction{-1}{\EXPfunction}{\Fexp}
	\PRODUCTfunction{\Fexp}{\SINfunction}{\FunTwo}
	
	 %%% 3. The Linear Combination
	 % Composing
	 \LINEARCOMBINATIONfunction{2}{\FunOne}
	 			{-1}{\FunTwo}{\LinearComb}
	  % Using
	 \LinearComb{4}{\FLinSol}{\DFLinSol}
	  % Printing
	 \begin{itemize}
	 	\item $F(4) = 2e^{cos 4} - e^{-4}sin 4$,
	 		which is \FLinSol
	 	\item $F'(t)= -2e^{cos t}sin t - e^{-t}cos t +
	 		e^{-t}sin t$, for $t=4$ we have \DFLinSol
	 \end{itemize}
	 
	 \pagebreak
	 
	 \subsection{Defining Polynomial Functions}
	 \begin{description}
	 	\item[Linear Polynomial] 
	 		$p(t) = a + b\,t$ \\
	 		\newlpoly{\linpoly}{4}{3} % a=4, b=3
	 		\linpoly{2}{\flinpoly}{\fdlinpoly}
	 		at $t=2$, it evaluates to $p(2)=\flinpoly$, 
			$p'(2)=\fdlinpoly$
		%
		\item[Quadratic Polynomial] 
			$p(t) = a + b\,t + c t^2$ \\
			\newqpoly{\quadpoly}{4}{3}{1} 
			\quadpoly{2}{\fquadpoly}{\fdquadpoly} 
			at $t=2$, it evaluates to $p(2)=\fquadpoly$, 
			$p'(2)=\fdquadpoly$
		%
		\item[Cubic Polynomial] 
			$p(t) = a + b\,t + c\,t^2 + d\,t^3$ \\
			\newcpoly{\cubepoly}{4}{3}{2}{1} 
			\cubepoly{2}{\fcubepoly}{\fdcubepoly} 
			at $t=2$, it evaluates to $p(2)=\fcubepoly$, 
			$p'(2)=\fdcubepoly$
	 \end{description}
	 
	\noindent
	Of course this whole setup can be abstracted inside 
	a new control sequence. \\[3pt]
	\newcommand{\quadpoly}[4]{
		% Define
		\newqpoly{\evquadpoly}{#1}{#2}{#3} % a,b,c
		% Compute
		\evquadpoly{#4}{\fquadpoly}{\fdquadpoly}
		% Print
		Evaluating $p(t)=a+bt+ct^2$ at $t=#4$ gives\,:-
		$p(#4)=\fquadpoly$, and $p'(#4)=\fdquadpoly$
	}
	\quadpoly{1}{2}{3}{4}
	
	\subsection{Vector Valued Functions}
	\[ 
		\mathbf{F}(t) = \begin{bmatrix} 
							t^2 \\ t^3 
						\end{bmatrix} 
	\]
	%% 1. V Define
	\VECTORfunction{\SQUAREfunction}{\CUBEfunction}{\VF}
	%% 2. V Compute
	\VF{2}{\fsq}{\fdsq}{\fcube}{\fdcube}
	%% 3. V Print
	If parameter/variable $t=2$ :
	\[ 
		F(2) = \begin{bmatrix} 
					\fsq   \rightarrow 2^2 \\ 
					\fcube \rightarrow 2^3 
			   \end{bmatrix} 
	\]
	and
	\[
		F'(2) = \begin{bmatrix} 
					\fdsq \\ 
					\fdcube
			   \end{bmatrix}	
	\]
	
	\subsection{Combining Two Polynomials}
	\[
		\mathbf{P}(t) = \begin{pmatrix}
							a + b\,t \\
							a + b\,t + c\,t^2
					    \end{pmatrix}			
	\]
	\newcommand{\PolVectFunc}[4]{ % a,b,c,t
		% Linear Polynomial
		\newlpoly{\vlinpoly}{#1}{#2}
		\vlinpoly{#4}{\fvlinpoly}{\fdvlinpoly}
		% Quadratic Polynomial
		\newqpoly{\vquadpoly}{#1}{#2}{#3}
		\vquadpoly{#4}{\fvquadpoly}{\fdvquadpoly}
		% Creating & Evaluating the vector function
		\VECTORfunction{\vlinpoly}{\vquadpoly}							{\VectFunc}
		\VectFunc{#4}{\fvlinpoly}{\fdvlinpoly}
			{\fvquadpoly}{\fdvquadpoly}
		%%
		Evaluating at $a=1,b=2,c=3, and~ t=4$, gives:
		\[
			P(4) = \begin{pmatrix}
						\fvlinpoly \\ \fvquadpoly
				   \end{pmatrix}			
		\]
		and
		\[
			P'(4) = \begin{pmatrix}
						\fdvlinpoly \\ \fdvquadpoly
				    \end{pmatrix}			
		\]
	}
	\PolVectFunc{10}{12}{35}{14}
	
	\subsubsection{Defining Custom Vector Functions}
	\[
		\mathbf{V}(t) = \begin{bmatrix}
							X(t) = t^2 \\ 
							Y(t) = t^3
						\end{bmatrix}
	\]
	%
	\newvectorfunction{\NewVF}{
		% Define (\x, \Dx) and (\y, \Dy) from \t
		\SQUARE{\t	}{\x}       % x  = t^2
		\MULTIPLY{2}{\t}{\Dx} % x' = 2t
		%
		\CUBE{\t}{\y}         % y = t^3
		\MULTIPLY{3}{\x}{\Dy} % y’= 3t^2
	}
	%
	\NewVF{2}{\XVF}{\XdVF}{\YVF}{\YdVF}
	\begin{center}
		$
			V(2) =  \begin{bmatrix}
						\XVF \\ \YVF
					\end{bmatrix}			
		$ ,
		$
			V'(2) =  \begin{bmatrix}
						\XdVF \\ \YdVF
					 \end{bmatrix}			
		$
	\end{center}
	
	\pagebreak
	
	\subsection{Defining Low Level New Functions}
	Defining $F(t) = t^2 + t^3$ which has a derivative 			of $F'(t) = 2t + 3t^2$ 
	%
	\newfunction{\myNewFunc}{
		\SQUARE{\t}{\a} % a = t^2
		\CUBE{\t}{\b}   % b = t^3
		\ADD{\a}{\b}{\y}% F(t) = t^2 + t^3
		%%
		\MULTIPLY{2}{\t}{\adash} % a' = 2t
		\MULTIPLY{3}{\a}{\bdash} % b' = 3*t^2 -> (a=t^2)
		\ADD{\adash}{\bdash}{\Dy}% F'(t) = 2t + 3t^2
	} 
	%
	\myNewFunc{2}{\myfsol}{\myfdsol}
	\[ F(t) = \myfsol,~ F'(t) = \myfdsol \]
\end{document}